\documentclass[11pt]{article}
\begin{document}
\section*{Step 40 Min-Cut Shadow-Price Statement}
This statement is scoped to the declared finite bulk graph carrier.

Let \(G\) be a finite capacitated bulk graph with boundary region \(A\) and complement \(\bar A\). The geometric area is computed as
\[
  \operatorname{Area}(A)=\min_{\gamma:A|\bar A}\sum_{e\in\gamma} c_e,
\]
the min-cut separating \(A\) from \(\bar A\). The entanglement-flow ledger is the primal max-flow from \(A\) to \(\bar A\),
\[
  S_{\rm flow}(A)=\max_f |f|,
\]
subject to the same edge-capacity constraints. Linear-program duality gives
\[
  S_{\rm flow}(A)=\operatorname{Area}(A)
\]
on the holographic graph. The boundary state used in the holographic carrier has reduced-density-matrix entropy equal to this flow capacity, so the computed coefficient
\[
  k = \operatorname{Area}(A)/S(A)
\]
is stable across the two refinements.

The non-holographic controls alter either the boundary state or the bulk capacity while keeping the other side fixed, and the equality fails. Thus the agreement is not a tautological rescaling of one parameter; it is a finite-carrier optimization-dual relation that recognition-lands on the Ryu--Takayanagi/bit-thread structure.
\end{document}
